\subsection{The five-object scope at eighteen arrows}

\begin{lemma}[Proved: semisimple corners in the five-object budget]
\label{CP2:A-five18-semisimple-small}
Let $C$ have five objects, acyclic support, eighteen arrows,
and corner orders at most two. Then $BC$ does not have the
homotopy type of $\CP^2$.
\end{lemma}
\begin{proof}
The rational algebra of every corner is semisimple. If the
support is not a full five-object chain, its strict height
is at most three. The relative bar resolution of
Lemma~\ref{CP2:A-semisimple-height} then has no homology in
degree four.

For the full chain there are ten strict hom-sets. Put
$D=\sum_x(|\operatorname{End}(x)|-1)$ and
$P=\sum_{x<y}(|C(x,y)|-1)$. The arrow count gives
$D+P=18-5-10=3$. If $D=0$, all corners are trivial and
at least one of the four consecutive hom-sets is a
singleton: otherwise $P\geq4$. Proposition~\ref{A7:adjacent-homs}
annihilates all degree-two products. If $D\geq1$, then
$P\leq2$, and the support cycle rank is $10-5+1=6$.
Lemma~\ref{CP2:A-high-two-parallel} applies; under the
assumption of projective-plane type its second Betti number
hypothesis is exactly one. Its zero-square conclusion again
contradicts the projective-plane cohomology ring.
\end{proof}

\begin{lemma}[Proved: deletion through a unique singleton comma]
\label{CP2:A-five18-singleton-deletion}
Suppose a category $C$ with acyclic support has an object $x$
such that its incoming arrows from other objects consist
of a single arrow $y\to x$, and $\operatorname{End}(y)$ is
trivial. The inclusion of the full deletion
$i:C\setminus\{x\}\longrightarrow C$ is a homotopy
equivalence on classifying spaces. The dual assertion
holds for a unique outgoing arrow to a trivial corner.
Consequently such a five-object eighteen-arrow category
does not have type $\CP^2$.
\end{lemma}
\begin{proof}
For every retained object $z$, the comma category
$(i\downarrow z)$ has terminal object $(z,1_z)$.
For $z=x$, that comma category consists of the one object
$(y,y\to x)$ and its identity, by the two stated hypotheses.
Quillen's Theorem~A therefore applies to this specified
inclusion. Opposition proves the dual assertion.
Deleting $x$ removes its identity and the specified cross
arrow, and possibly more arrows, so the deletion has four
objects and at most sixteen arrows. The independently
proved complete sixteen-arrow exclusion rules it out.
\end{proof}

\begin{lemma}[Proved: the tree-edge scope test]
\label{CP2:A-five18-tree-edges}
Let $C$ have connected acyclic support whose undirected
comparability graph is a tree. Suppose every nonleaf corner
has order at most five, and every edge with endpoint corner
orders $a,b$ and stalk size $s$ satisfies at least one of
\[
\begin{gathered}
s\leq2,\qquad \max(a,b)\leq4,\qquad a+b+s\leq11,\\
\min(a,b)\leq2\ \text{ and }\ \max(a,b)\leq11,
\qquad \{a,b\}=\{3,5\}.
\end{gathered}
\]
If $BC$ has bounded rational homology, all its rational
degree-two products vanish.
\end{lemma}
\begin{proof}
The tree condition forbids a strict three-object chain.
Thus the nerve is the union of its full edge subcategories,
attached successively along corner categories. Every
intersection corner is a nonleaf corner. By
Corollary~\ref{CP2:A-five-types}, its rational homology
vanishes in degrees at least three. The Mayer--Vietoris
sequences consequently give injectivity of restriction in
degree four, and show that every edge piece inherits
bounded rational homology. These statements also apply
when the intersection corner has the type of $S^2$;
rational acyclicity of that corner is not required.

For $s=1$ the edge space is the join of its two corner
spaces and has zero second rational cohomology. For $s=2$
use Lemma~\ref{CP2:A-high-two-cross}. When both corners
have order at most four, both restrictions in degree two
are zero, so Lemma~\ref{CP2:A-two-block} applies. For
$a+b+s\leq11$ use Corollary~\ref{CP2:A-high-two-objects}.

If one corner has order at most two, its trivial rational
modules are projective on both sides. The appropriate
actual orbit functor of Lemma~\ref{CP2:A-high-orbits}
is a rational cohomology-ring equivalence to an ordinary
act cone over the other corner. Boundedness passes through
this equivalence, so Theorem~\ref{CP2:A-bounded-eleven-acts}
applies when the other corner has order at most eleven.

Finally consider orders three and five. If the five-element
corner is not the rectangular-band monoid $P_5$, both
corners have zero rational second cohomology and the
two-block lemma applies. Otherwise the three-element
corner either has the required projective trivial module,
giving an ordinary $P_5$ act cone by the same orbit
comparison, or is the exceptional oriented right-zero
band monoid. The ordinary act is covered by
Lemma~\ref{CP2:A-high-two-row-act}, and the exceptional
biset by Lemma~\ref{CP2:A-high-P5-graph}. These arguments
allow every nonempty finite stalk. In all five cases the
edge products vanish. Injectivity of fourth restriction
therefore kills the global products.
\end{proof}

\paragraph{Complete interpretation of the cardinality predicates.}
The source \path{code/five18_scope.py} considers only connected
acyclic supports on exactly five objects and positive
decorations of total eighteen. For each support its cycle
rank is at most six. Its support-first canonization chooses
one representative under object relabelling and opposition,
and then divides decorations by the complete automorphism
and anti-automorphism group of that fixed support.
No corner multiplication or composition table is assumed
to exist merely because its cardinalities have been listed.

Here is the exact mathematical meaning of every enabled
proof predicate. Nonzero first mod-two homology of the
support is excluded by the fundamental-group surjection:
lift each traversed support edge to any category arrow,
using the reversed path when necessary. Singleton cross
hom-sets use Lemma~\ref{CP2:A-singleton-support}; every
second class pulls back from the support, whose fourth
cohomology is zero by Lemma~\ref{CP2:A-high-flag}.
Corners at most two use
Lemma~\ref{CP2:A-five18-semisimple-small}. Parallel excess
at most two with corners at most four uses
Lemma~\ref{CP2:A-high-two-parallel} with the candidate's
second Betti number one. Nonserial parallel hom-sets use
Lemma~\ref{CP2:A-high-parallel-cut}: a small corner has
zero second cohomology, and a corner incident with a
singleton cross hom-set has nullhomotopic inclusion into
that edge join, so every global second class restricts
to zero there.

For every proper downset and complementary upset, the
source tests the two-block lemma and the fixed-arrow
corollary. Its sufficient zero-second-cohomology test is
applied component by component. A one-object component
has corner order at most four. A larger component either
has at most $2k+1$ arrows, covered by
Lemma~\ref{CP2:A-high-H2-cost}; has singleton cross
hom-sets, covered by the singleton-support lemma and the
flag budget on at most four vertices; or has semisimple
corners of orders at most two and strict height at most
one. This last case follows from the relative bar bound.
These are sufficient tests only. A failed test removes
nothing. A full singleton join across a cut has zero
positive products, since the join is a suspension after
collapsing its two contractible basepoint cones. A small
extreme corner with every incident stalk nonempty and of
size at most two uses Lemma~\ref{CP2:A-high-small-binary},
again with the candidate's second Betti number one.

Unique singleton comma deletions use
Lemma~\ref{CP2:A-five18-singleton-deletion}. The tree-edge
predicate is exactly Lemma~\ref{CP2:A-five18-tree-edges}.
A full star with center at most eleven and every other
corner at most two is a special case of
Proposition~\ref{CP2:A-tree-one-large-corner-eleven}.
Thus every recorded proof reason is a sufficient exclusion
under explicitly stated candidate hypotheses. The numerical
profile and residual joins are separate finite obligations.

\paragraph{Reproducible orbit counts and the separate tree supplement.}
For a fixed support with $q$ nonempty hom-sets, including
the five diagonal positions, all positive decorations
number $\binom{17}{q-1}$. The source lists every positive
composition of eighteen into these $q$ positions, retains
the least decoration in each support symmetry orbit, and
records its orbit size as group order divided by its
stabilizer order. The orbit-size sum is checked against
that binomial coefficient for every support. The full
stream preserves each profile and every applicable named
reason. An independent program must still regenerate the
full list and its predicate partition before the numerical
scope is adopted in a completed lower bound.

The later source \path{code/five18_tree_supplement.py}
leaves this first stream and its residual ledger unchanged.
It checks the exact hypothesis of
Proposition~\ref{CP2:A-tree-one-large-corner-eleven} on
each residual row and supplies a separately hashed proof
partition. Rows not covered by it remain finite-composition
obligations; no running or incomplete search excludes one.
